\documentclass[a4,11pt]{aleph-notas}
% Se puede ver la documentación aquí:
% https://github.com/alephsub0/LaTeX_aleph-notas

% -- Paquetes adicionales
\usepackage{enumitem}
\usepackage{url}
\usepackage{array}
\usepackage{booktabs}
\usepackage{longtable}
\hypersetup{
    urlcolor=blue,
    linkcolor=blue,
}

% -- Datos
\institucion{Facultad de Ciencias Exactas, Naturales y Ambientales}
\carrera{Catálogo STEM}
\asignatura{Álgebra lineal}
\tema{Clase 05: El espacio $\mathbb{R}^n$}
\autor{Andrés Merino}
\fecha{Periodo 2026-01}

\logouno[0.16\textwidth]{Logos/logoPUCE_04_ac}
\definecolor{colortext}{HTML}{1957d1}
\definecolor{colordef}{HTML}{1957d1}
\fuente{montserrat}

\begin{document}

\encabezado 

%%%%%%%%%%%%%%%%%%%%%%%%%%%%%%%%%%%%%%%%
\section*{Resultado de Aprendizaje}
%%%%%%%%%%%%%%%%%%%%%%%%%%%%%%%%%%%%%%%%

%%%%%%%%%%%%%%%%%%%%%%%%%%%%%%%%%%%%%%%%
\subsection*{RdA de la asignatura:}
\begin{itemize}[leftmargin=*]
    \item \textbf{RdA 1:} Comprender los conceptos fundamentales del Álgebra Lineal, incluyendo el estudio de matrices, determinantes, sistemas de ecuaciones lineales y espacios vectoriales, destacando su importancia en el análisis y resolución de problemas matemáticos.
    \item \textbf{RdA 2:} Resolver problemas prácticos y teóricos relacionados con el Álgebra Lineal, utilizando técnicas de cálculo con matrices, determinantes y sistemas de ecuaciones, así como en el análisis de espacios vectoriales.
\end{itemize}


%%%%%%%%%%%%%%%%%%%%%%%%%%%%%%%%%%%%%%%%
\subsection*{RdA de la actividad:}
\begin{itemize}[leftmargin=*]
    \item Comprender la estructura de $\mathbb{R}^n$ y su base canónica.
    \item Aplicar producto punto, norma y distancia para resolver ejercicios en $\mathbb{R}^n$.
    \item Interpretar relaciones geométricas entre vectores: ortogonalidad, paralelismo y ángulo.
    \item Modelar rectas y planos mediante ecuaciones paramétricas y cartesianas en $\mathbb{R}^2$ y $\mathbb{R}^3$.
\end{itemize}

%%%%%%%%%%%%%%%%%%%%%%%%%%%%%%%%%%%%%%%%
\section*{Introducción}
%%%%%%%%%%%%%%%%%%%%%%%%%%%%%%%%%%%%%%%%

%%%%%%%%%%%%%%%%%%%%%%%%%%%%%%%%%%%%%%%%
\paragraph{Control de lectura:}
Antes de iniciar la clase, se aplicará el control de lectura del siguiente enlace:
\href{https://gemini.google.com/share/240a32bed915}{Control de lectura - Clase 05}.

\paragraph{Pregunta inicial:}
¿Cómo podemos trabajar con vectores de más de dos o tres dimensiones? ¿Podemos hacer geometría en $\mathbb{R}^5$?

%%%%%%%%%%%%%%%%%%%%%%%%%%%%%%%%%%%%%%%%
\section*{Desarrollo}
%%%%%%%%%%%%%%%%%%%%%%%%%%%%%%%%%%%%%%%%

%%%%%%%%%%%%%%%%%%%%%%%%%%%%%%%%%%%%%%%%
\subsection*{Actividad 1: Operaciones y propiedades en $\mathbb{R}^n$}

La clase inicia con recapitulación de la idea intuitiva de vectores de 2 y 3 dimensiones. Luego, mediante clase magistral y práctica guiada, se formalizan operaciones con vectores, producto punto, norma y distancia para resolver problemas con interpretación geométrica.

\paragraph{¿Cómo lo haremos?}
\begin{itemize}[leftmargin=*]
    \item \textbf{Clase magistral:} se explican producto punto, norma, distancia, ángulo, ortogonalidad y proyección. Se usa el resumen \href{https://fcena-puce.github.io/AlgebraLineal-05-N0048/01Resumen/Resumen05.pdf}{Resumen05.pdf}.
    \item \textbf{Resolución de ejercicios:} se guiará a los estudiantes en la solución de ejercicios del documento \href{https://fcena-puce.github.io/AlgebraLineal-05-N0048/04Ejercicios/Ejercicios05.pdf}{Ejercicios05.pdf}, aplicando los conceptos vistos.
    \item \textbf{Aplicación geométrica:} se construyen ecuaciones paramétricas y cartesianas de rectas en $\mathbb{R}^2$ y planos en $\mathbb{R}^3$, visualizando en Geogebra.
\end{itemize}

\paragraph{Verificación de aprendizaje:}
\begin{enumerate}[leftmargin=*]
    \item ¿Cómo se calcula el producto punto y qué información geométrica aporta?
    \item ¿Cuándo dos vectores son ortogonales o paralelos en $\mathbb{R}^n$?
    \item ¿Cómo se obtiene la ecuación cartesiana a partir de una ecuación paramétrica de recta o plano?
\end{enumerate}

%%%%%%%%%%%%%%%%%%%%%%%%%%%%%%%%%%%%%%%%
\section*{Cierre}
%%%%%%%%%%%%%%%%%%%%%%%%%%%%%%%%%%%%%%%%

\paragraph{Tarea:}
Resolver del libro \href{https://catalogobiblioteca.puce.edu.ec/cgi-bin/koha/opac-detail.pl?biblionumber=86081}{Fundamentos de álgebra lineal de Larson}, sección 4.1, los ejercicios: 7, 11, 19, 23, 31, 33, 39, 41.

\paragraph{Pregunta de investigación:}
\begin{enumerate}[leftmargin=*]
    \item ¿Por qué la desigualdad de Cauchy-Schwarz garantiza que el ángulo entre vectores está bien definido?
    \item ¿Qué interpretación geométrica tiene la proyección ortogonal en problemas aplicados?
    \item ¿Cómo se puede decidir si dos rectas en $\mathbb{R}^3$ se intersectan, son paralelas o se cruzan usando su forma paramétrica?
\end{enumerate}

\paragraph{Para la próxima clase:}  
Visualizar el video \href{https://youtu.be/fNk_zzaMoSs?si=o7DzunwZHW_YP4lg}{Vectores | Capítulo 1, Esencia del álgebra lineal}. Luego, leer la sección 4.2 del libro \href{https://catalogobiblioteca.puce.edu.ec/cgi-bin/koha/opac-detail.pl?biblionumber=86081}{Fundamentos de álgebra lineal de Larson}.




\end{document}
